\documentclass[../main.tex]{subfiles}
\begin{document}

Chương này giới thiệu bài toán mà luận văn giải quyết, xuất phát từ nhu cầu cải tiến liên tục các khoá đào tạo nội bộ trong doanh nghiệp phần mềm. Phần \ref{sec:dvd} trình bày lý do bài toán này cần một hướng tiếp cận mới. Phần \ref{sec:giaiphap} tổng hợp các hướng nghiên cứu liên quan và chỉ ra khoảng trống mà các kiến trúc hiện có chưa giải quyết. Từ khoảng trống đó, phần tiếp theo trình bày mục tiêu nghiên cứu, khung khái niệm và đóng góp chính của luận văn, trước khi khép lại bằng bố cục của toàn bộ báo cáo.

\section{Problem Statement}
\label{sec:dvd}

Trong bối cảnh corporate training tại các doanh nghiệp phần mềm, chương trình đào tạo nội bộ hiếm khi là một sản phẩm thiết kế một lần. Các khoá học dành cho nhân viên mới, intern, fresher hoặc kỹ sư chuyển đổi vai trò thường phải được cập nhật liên tục theo thay đổi của công nghệ, yêu cầu dự án, đặc điểm người học và phản hồi sau mỗi đợt triển khai. Trong thực tế, một module về Java, lập trình hướng đối tượng, cấu trúc dữ liệu, kiểm thử phần mềm hoặc quy trình phát triển phần mềm có thể cần được cải tiến nhiều lần về mục tiêu học tập, bài tập, rubric, tài liệu giảng dạy, ràng buộc triển khai và kế hoạch đánh giá.

Vấn đề này xuất phát từ trải nghiệm thực tế trong quá trình xây dựng và cải tiến các khoá học nội bộ cho nhân viên, intern và fresher. Trong luận văn này, instructional design (thiết kế dạy học) được hiểu là quá trình thiết kế có hệ thống các trải nghiệm học tập, bao gồm phân tích nhu cầu và người học, xác định mục tiêu, thiết kế đánh giá, phát triển hoạt động và tài liệu, lập kế hoạch triển khai, và cải tiến dựa trên phản hồi. Các \gls{LLM} đủ mạnh hiện nay đã có thể đảm đương gần như toàn bộ quá trình này, chứ không chỉ hỗ trợ từng tác vụ riêng lẻ. Tuy nhiên, khi được sử dụng như một chatbot độc lập, chúng còn hạn chế trong việc duy trì một vòng đời \gls{ISD} nhất quán cho toàn bộ khoá học: mỗi lượt tương tác thường xử lý độc lập một tác vụ mà không giữ được trạng thái ràng buộc giữa các thành phần thiết kế đã sinh trước đó. Đồng thời, khi năng lực model nền dần bão hoà trên nhiều benchmark \cite{deveci2025ouroboros,shukla2025scalingwall}, khác biệt về chất lượng đầu ra của các hệ agent ngày càng phụ thuộc vào lớp harness bao quanh model hơn là bản thân model \cite{cho2026harness}; lớp harness này càng có giá trị khi doanh nghiệp chủ động dùng \gls{SLM} để tối ưu chi phí vận hành, vì model càng nhỏ càng cần cơ chế kiểm tra, sửa lỗi và tinh chỉnh bên ngoài bù lại.

Bài toán đặt ra cho luận văn, do đó, không chỉ là sinh nội dung ISD một lần cho đủ, mà là xây dựng một kiến trúc agent có khả năng biểu diễn, kiểm tra và sửa lại các thành phần thiết kế dạy học một cách có mục tiêu, mỗi khi một phần của khoá học cần được cải tiến.

\section{Background and Problems of Research}
\label{sec:giaiphap}

Các phương pháp gần đây trong lĩnh vực AI for education đã tiếp cận bài toán hỗ trợ thiết kế dạy học từ nhiều hướng khác nhau. Một nhóm nghiên cứu tập trung xây dựng benchmark để đánh giá năng lực giáo dục của LLM, tiêu biểu là OpenLearnLM, EduBench, MathTutorBench và TutorBench, nhưng các benchmark này chủ yếu đánh giá kiến thức, kỹ năng hoặc năng lực tutoring của mô hình, thay vì đánh giá toàn bộ vòng đời ISD của một khoá học. Một nhóm khác tập trung vào hỗ trợ thiết kế bài học hoặc lập kế hoạch dạy học, như LearnLM, LessonPlanner, EduPlanner và các nghiên cứu về sinh giáo án (teaching plan generation) bằng GPT-4; các hệ thống này cải thiện chất lượng của một artifact đơn lẻ nhưng chưa đặt artifact đó trong một quy trình ISD có kiểm soát chất lượng xuyên suốt. Gần với bài toán của luận văn hơn cả là ISD-Agent-Bench, benchmark đề xuất dựa trên mô hình ADDIE cùng các agent dựa trên lý thuyết như ReAct-ADDIE, ADDIE-Agent, Dick-Carey-Agent và RPISD-Agent. Các hướng nghiên cứu này minh hoạ nhiều cách kiểm soát quá trình sinh artifact khác nhau, từ one-shot prompting, tool-use, workflow tuần tự, đến multi-agent generation và các mô hình ISD cổ điển; phân tích chi tiết về kiến trúc, điểm mạnh và giới hạn của từng hướng được trình bày trong Chương 2.

Tuy vậy, các kiến trúc hiện có vẫn chủ yếu tối ưu cho chất lượng artifact tương thích ADDIE tại lần sinh đầu tiên. Chúng chưa đặt mức độ alignment giữa các thành phần thiết kế như một trạng thái trung tâm có thể truy vết, kiểm tra, sửa chữa hay tinh chỉnh một cách hiệu quả. Đây là một giới hạn quan trọng đối với bài toán cải tiến khoá học liên tục, vì chất lượng thiết kế không chỉ phụ thuộc vào việc có đủ mục tiêu, hoạt động, bài tập, rubric và kế hoạch triển khai, mà còn phụ thuộc vào việc các thành phần đó có nhất quán với nhau hay không. Nói cách khác, hệ thống cần biết mục tiêu nào được đo bởi bài đánh giá nào, hoạt động nào hỗ trợ mục tiêu nào, rubric nào phản ánh đúng evidence of learning nào, và ràng buộc triển khai nào có thể làm kế hoạch dạy học không khả thi. Khi một thành phần bị thay đổi để sửa lỗi hoặc cải tiến, các kiến trúc tuyến tính hiện có không có cơ chế tường minh để lan truyền ảnh hưởng đó tới các thành phần liên quan, khiến mỗi lượt cải tiến thường phải rà soát hoặc sinh lại trên diện rộng — khó kiểm soát và kém hiệu quả so với nhu cầu cải tiến liên tục.

\section{Research Objectives and Conceptual Framework}

Từ khoảng trống nêu trên, luận văn đặt mục tiêu xây dựng một kiến trúc agent trong đó mức độ alignment giữa các thành phần thiết kế dạy học được biểu diễn tường minh, thay vì chỉ ẩn trong quá trình sinh artifact. Luận văn đề xuất \textbf{AlignmentGraph-ISD}, một alignment-graph harness cho ISD agent, theo công thức tổng quát:

\begin{equation}\label{eq:formula}
    \text{AlignmentGraph-ISD Agent} = \text{AlignmentGraph-ISD Harness} + \text{Model}
\end{equation}

Công thức~\eqref{eq:formula} thể hiện nguyên tắc thiết kế cốt lõi: harness cung cấp môi trường thực thi chuyên cho ISD, còn model cung cấp năng lực ngôn ngữ, suy luận và sinh nội dung bên trong môi trường đó. Harness không gắn cứng với một mô hình thiết kế dạy học duy nhất; thay vào đó, nó biểu diễn ISD workflow như một graph có thể cấu hình, trong đó ADDIE đóng vai trò output contract chung để tương thích với ISD-Agent-Bench, còn các chiến lược như backward design, constructive alignment, Gagne sequencing hoặc rubric-first design được dùng như policy và verifier bên trong. Khung khái niệm này được cụ thể hoá trong Chương 3 thành năm lớp kiến trúc: từ scenario hoặc yêu cầu giảng viên, qua orchestrator và dynamic ISD task graph, qua four-layer memory và tool layer, qua schema validation/pedagogical verification/targeted repair, đến ADDIE output aggregator và trajectory.

Để đánh giá kiến trúc đề xuất, luận văn đặt ra hai câu hỏi nghiên cứu. \textbf{RQ1}: Trên ISD-Agent-Bench, AlignmentGraph-ISD cải thiện chất lượng artifact instructional design như thế nào so với các agent baseline và các ISD agent hiện có? \textbf{RQ2}: AlignmentGraph-ISD có hỗ trợ tinh chỉnh khoá học qua nhiều phiên một cách hiệu quả không — một yêu cầu chỉnh sửa có được khoanh vùng ảnh hưởng, chỉ sửa đúng các thành phần liên quan với chi phí thấp hơn sinh lại toàn bộ, mà không làm suy giảm alignment? RQ1 được đo định lượng trên benchmark (Chương 5), với mỗi scenario là một lượt thiết kế đầu tiên; RQ2 được đo bằng các process metric trên các phiên tinh chỉnh của case study, theo giao thức định nghĩa ở Chương 4. \gls{RAG} trên external corpus cũng được đánh giá riêng qua case study, vì đây là một add-on chỉ dùng cho ứng dụng thực tế và không nằm trong phạm vi benchmark.

\section{Contributions}

Luận văn đề xuất một đóng góp chính: \textbf{AlignmentGraph-ISD}, một alignment-graph harness cho \gls{ISD} agent. Harness này đặt model trong một execution framework chuyên cho ISD, gồm dynamic task graph, multi-agent orchestration, four-layer memory, graph-context retrieval (core) và agentic RAG trên external corpus (case study), schema-constrained generation, pedagogical verification, targeted repair và workflow trace. Thay vì chỉ sinh một artifact cuối, AlignmentGraph-ISD tổ chức quá trình thiết kế dạy học thành các artifact node và alignment edge có quan hệ rõ ràng, cho phép kiểm soát chất lượng, tái chạy phần lỗi có mục tiêu và thu thập dữ liệu đánh giá quá trình.

Bên cạnh đóng góp chính, luận văn trình bày một case study trong miền IT/Software Engineering, minh hoạ cách AlignmentGraph-ISD được dùng trong một workflow thực tế: nhận yêu cầu thiết kế hoặc cải tiến một module đào tạo, truy xuất tài liệu miền phần mềm, tạo ADDIE-compatible artifact, kiểm tra alignment giữa mục tiêu, hoạt động và đánh giá, sinh rubric và bài tập, ghi lại trajectory, và xuất kết quả theo định dạng có thể đưa vào pipeline đánh giá. Case study không nhằm chứng minh hiệu quả học tập dài hạn của học viên, mà nhằm cho thấy AlignmentGraph-ISD có thể trở thành một công cụ hỗ trợ instructional designer hoặc giảng viên nội bộ trong bối cảnh đào tạo phần mềm.

\section{Organization of Thesis}

Phần còn lại của báo cáo được tổ chức như sau. Chương 2 trình bày literature review, tập trung vào các kiến trúc agent và benchmark liên quan trong lĩnh vực AI/LLM trước khi tóm lược các khái niệm ISD nền tảng cần thiết để đọc phần phương pháp. Chương 3 trình bày phương pháp đề xuất, gồm mô hình alignment graph, workflow động, cơ chế multi-agent orchestration, schema validation, pedagogical verification và four-layer memory của harness. Chương 4 trình bày case study trong miền IT/Software Engineering, mô tả kiến trúc ứng dụng tích hợp harness với quản lý phiên nhiều lượt, cơ chế truy xuất tri thức theo domain, giao diện prototype, và giao thức đo hiệu quả tinh chỉnh cho RQ2. Chương 5 trình bày kết quả đánh giá cho cả hai câu hỏi nghiên cứu: kết quả định lượng trên ISD-Agent-Bench gồm kết quả sơ bộ trên subset \texttt{test\_30} và kết quả mở rộng trên subset \texttt{test\_90}, hiện trạng đo lường cho RQ2, cùng thiết kế ablation bổ trợ và các threats to validity. Chương 6 tổng kết các kết quả đạt được và đề xuất hướng phát triển tiếp theo.

\end{document}
